% Ad Familiares 4.5 --- headnote
% ad-familiares-04-05, mid-March 45 BC, Athens

\textit{Servius Sulpicius Rufus to Cicero, written at Athens
about the middle of March 45 BC (Perseus: \textit{Athenis
circ.~medio m.~Mart.~a.~709 (45)}). Tullia, Cicero's daughter,
had died at Tusculum about a month earlier, in mid-February;
Cicero's letter informing his old friend would have reached
Servius in Greece a few weeks before this reply was sent back.
Servius, a great jurist and a man of austere Stoic temperament,
was at the time serving as proconsul of Achaea.}

\textit{The letter is one of the small masterpieces of the Roman
\textit{consolatio}, and it is meant to be read together with
Cicero's reply (\textit{Fam.}~4.6). What Servius writes is what
he believes Cicero needs to hear rather than what would feel
kindest: a sustained argument that Tullia's death is no greater
calamity than the wider losses already borne, that her continued
life at this hour had little to offer her, and that the proper
office of the philosophical man is to apply to himself the
prescriptions he has given to others. The argument turns on the
famous passage in \S 4: returning from Asia and sailing from
Aegina toward Megara, Servius looked round at the ruined
\textit{oppida cadavera} of Aegina, Megara, the Piraeus, and
Corinth, and saw the corpses of cities laid out in one place ---
and asked himself by what right ``us little men''
(\textit{nos homunculi}), whose lives were always the briefer
thing, take the death of one of our own as an outrage.}

\textit{The voice across the letter is consistent: a
jurist-philosopher's voice, drier and more disciplined than
Cicero's own, with the Stoic premise that grief beyond a certain
measure is itself a failure of \textit{ratio}. The closing line
of \S 6 --- that out of all the virtues, this one alone, the
bearing of adversity, should not seem to be missing in Cicero ---
is offered as a challenge as much as a consolation. Cicero's
reply, written a few weeks later at Atticus's Ficuleanum estate,
accepts the philosophical argument but reports, with a candour
that goes against everything Servius has said, that the
consolation has not yet reached the wound.}
